(1) 热机连接恒温热源Th和冷源Tc，Th>Tc>0。热机内部可逆地在温度T1、T2运行，其中Th>T1>T2>Tc；两侧线性热导均为K，热流Qh=K(Th-T1)、Qc=K(T2-Tc)。写出内可逆约束与热机功率。
(2) 在T1,T2可调但保持各温差为正的条件下，严格求最大输出功率的内部温度和效率。
(3) 计算两个传热接触界面的总熵产生率，证明其非负，解释为什么机器内部可逆而整体不可逆。
(4) 求K趋零、Th趋Tc及热机内部温差消失时的极限；解释所得效率与Carnot效率的差别。
